۱۹۶  
داک وقتی که مستری از روی هوس برای یک سمینار دوست‌یابی که داک در یادگیری آنکس برگزار می‌کرد، ثبت‌نام کرد، با جامعه آشنا شد. مستری به‌طور صبورانه گوش داد در حالی‌که داک نکات و تاکتیک‌هایی را به اشتراک می‌گذاشت که نسبت به تکنولوژی موجود در جامعه مواردی مقدماتی AFC بودند. بعد از آن، او با داک صحبت کرد که اقرار کرد چندان مردی برای خانم‌ها نیست. بنابراین مستری او را برای یک شب بیرون برد، او را در روش مستری آموزش داد و به او امکان دسترسی به لانژ را داد. حالا داک یک ماشین شده بود، با حرمسرای خودش از زن‌ها. نام مستعار او از دکترایش در روانشناسی آمده بود، بنابراین با او تماس گرفتم و از او مشورت خواستم.

او پیشنهاد کرد که این سوالات را از مستری دقیقاً به این ترتیب بپرسم:
- آیا آن‌قدر افسرده‌ای که حس کنی می‌خواهی از همه چیز دست بکشی؟
- آیا زیاد درباره مرگ فکر می‌کنی؟
- آیا به فکر آسیب رساندن به خودت یا انجام کاری مخرب هستی؟
- آیا به خودکشی فکر می‌کنی؟
- چگونه این کار را انجام می‌دهی؟
- چه چیزی تو را از انجام دادن آن باز می‌دارد؟
- آیا فکر می‌کنی که این کار را در ۲۴ ساعت آینده انجام می‌دهی؟

من سوال‌ها را روی یک برگ کاغذ نوشتم، آن را به چهار تا تا زدم و در جیب پشتی‌ام قرار دادم. این یواشکی‌ام می‌بود. روال کاری‌ام.

وقتی به خانه مستری رسیدم، او در حال جداسازی تخت‌اش بود. حرکات او مکانیکی بود. پاسخ‌های او هم همین‌طور.

\- سبک: داری چی کار می‌کنی؟

\- مستری: دارم تخت‌ام را به خواهرم می‌دهم. او را دوست دارم و او لایق تخت بهتری است.

\- سبک: آیا آن‌قدر افسرده‌ای که حس کنی می‌خواهی از همه چیز دست بکشی؟

\- مستری: بله. بی‌فایدگی‌اش است. این میمتیک است. اگر میمتیک‌ها را درک کنید، آنگاه می‌فهمید که همه‌اش بی‌فایده است. نکته‌ای نیست.

\- سبک: اما تو یک ذهن برتر داری. این وظیفه توست که بقا ایجاد کنی.

\- مستری: مهم نیست. می‌خواهم ژن‌هایم را از بین ببرم.

\- سبک: آیا زیاد درباره مرگ فکر می‌کنی؟

\- مستری: همیشه.

\- سبک: آیا به فکر آسیب رساندن به خودت یا انجام کاری مخرب هستی؟

\- مستری: بله. این زندگی فابار است.

\- سبک: آیا به خودکشی فکر می‌کنی؟

\- مستری: بله.

\- سبک: چگونه این کار را انجام می‌دهی؟

\- مستری: غرق شدن، چون از آن بیشتر می‌ترسم.

\- سبک: چه چیزی تو را از انجام دادن آن باز می‌دارد؟

\- مستری: باید همه وسایلم را ببخشم. کامپیوتر پاتریشیا را انداختم و شکستم. بنابراین می‌خواهم کامپیوترم را به او بدهم. او نیاز به کامپیوتر دارد.

۱۹ب  
\- سبک: آیا او اهمیتی می‌داد؟

\- مستری: نه، واقعاً نه.

\- سبک: آیا او از این‌که آن را شکستی عصبانی بود؟

\- مستری: نه.

\- سبک: آیا فکر می‌کنی که زندگی‌ات را در ۲۴ ساعت آینده می‌گیری؟

من مستری را در اتاقش تنها گذاشتم، به آشپزخانه رفتم و شماره اطلاعات را برای شماره والدینش گرفتم. نام واقعی او اریک فون مارکوویک بود، اما آن هم یک توهم دیگر بود. آن را به‌طور قانونی از نام تولدش، اریک هورووات-مارکوویک، تغییر داده بود.

تلفن یک بار، دو بار، سه بار زنگ خورد. مردی گوشی را برداشت. صدایش خشن بود، طرز برخوردش تند. او پدر مستری بود.

\- [زنگ در به صدا در می‌آید]

"سلام، من دوست پسر شما، اریک، هستم."

\- سبک: کیه؟

"کی هستی؟"

\- "من نیل هستم، دوست اریک. و می‌خواستم..."

"دیگر اینجا تماس نگیر!" او عصبانی گفت.

\- صدای اینترکام: سلام، اینجا تایلر داردن است. من اینجا برای مستری هستم. از پست‌های او خوشم می‌آید و می‌خواهم ببینم می‌توانم او را ملاقات کنم.

"اما او نیاز..."

\- سبک: احتمالاً الان وقت خوبی نیست. صدای قطع شدن.

فقط یک نفر دیگر بود که می‌توانستم تماس بگیرم. به اتاق مستری برگشتم. او در حال بلعیدن قرصی با یک لیوان آب بود. صورتش قرمز و پیچیده بود، انگار که اشک‌های نامرئی می‌ریخت.

\- صدای اینترکام: اما من از کینگستون آمده‌ام.

\- سبک: ببخشید، مرد. او نمی‌تواند کسی را ببیند. او، ام، مریض است.

"چی خوردی؟" پرسیدم.

"چند قرص خواب‌آور"، گفت.

"چند تا؟" لعنتی. باید آمبولانس می‌زدم.

"دو."

"چرا خوردی؟"

"وقتی بیدارم، زندگی افتضاح است. بی‌فایده است. وقتی خوابم، خواب می‌بینم." او شروع به صحبت کردن مثل مارلون براندو در آخرالزمان حالا کرد. "دیشب خواب دیدم که در یک دلورین پرنده هستم. مثل آن چیزی که در بازگشت به آینده بود. و همه این سیم‌ها دور ما بود. من با خواهرم بودم. و او در حال رانندگی بود. ما بالای سیم‌ها رفتیم. و زندگی‌ام را زیر آن‌ها دیدم."

"گوش کن"، گفتم. "شماره تلفن پاتریشیا را لازم دارم."

حالا اشک‌هایش آمد. او مثل یک بچه بزرگ به‌نظر می‌رسید. یک بچه بزرگ که می‌خواست خودش را بکشد.

"می‌توانی شماره تلفن پاتریشیا را به من بدهی؟" دوباره به‌آرامی و ملایمت پرسیدم، طوری‌که انگار با یک کودک صحبت می‌کردم.

او شماره‌اش را به من داد – آهسته و به آرامی، مانند یک کودک.